\chapter{Percutaneous Lung Biopsy}

\section*{Overview}
Percutaneous lung biopsy is a procedure that uses a needle through the chest wall to sample a lung nodule or mass for diagnosis.

\section*{Alternative Names}
Transthoracic needle biopsy; CT-guided lung biopsy; lung needle aspiration

\section*{Principle}
Under CT or ultrasound guidance, a needle is advanced into the lesion and tissue or cells are aspirated for cytology and histology.

\section*{History}
Transthoracic needle biopsy developed in the mid-20th century and became a standard, image-guided method for diagnosing peripheral lung lesions.

\section*{Why It Is Done}
Performed for peripheral lung nodules or masses when bronchoscopy cannot reach them and malignancy is suspected.

\section*{Is This a Primary or Secondary Test or Procedure}
Primary diagnostic test for selected peripheral lung lesions.

\section*{How It Is Performed}
Performed under local anesthesia with CT guidance. A needle is advanced into the lesion, samples are taken, and the patient is positioned to minimize pneumothorax risk.

\section*{Preparation}
Imaging review, and holding anticoagulants.

\section*{Contraindications and Precautions}
Severe emphysema, bleeding risk, and inability to tolerate pneumothorax are relative contraindications.

\section*{Risks}
Pneumothorax, bleeding, hemoptysis, and infection.

\section*{Aftercare and Recovery}
A chest X-ray is often taken to check for pneumothorax.

\section*{Outcome and Follow-up}
Cytology and histology determine benign versus malignant disease.

\section*{Expected Results}
Benign tissue with no evidence of malignancy.

\section*{Follow-up}
Repeat imaging or surgical biopsy if sampling is nondiagnostic.

\section*{References}
\begin{enumerate}
  \item MedlinePlus. Lung biopsy. U.S. National Library of Medicine; 2025.
  \item Current Medical Diagnosis and Treatment. McGraw-Hill; 2025.
\end{enumerate}

\clearpage
